% --------------------------------------------------------------
% This is all preamble stuff that you don't have to worry about.
% Head down to where it says "Start here"
% --------------------------------------------------------------
 
\documentclass[12pt]{article}

\usepackage[spanish]{babel}
\usepackage[a4paper, margin=1in, textwidth=7in]{geometry} % Adjusting the textwidth and margins
 
\usepackage[margin=1in]{geometry} 
\usepackage{amsmath,amsthm,amssymb,siunitx,graphicx,icomma,parskip,mathtools,placeins,cancel}
\usepackage[hidelinks]{hyperref}
 
% Define un contador para los enunciados
\newcounter{pregun}
\newcounter{ejer}

% Define el entorno pregunta
\newenvironment{pregunta}[1][Cuestión]{%
    \refstepcounter{pregun} % Incrementa el contador y lo hace referenciable
    \begin{quote} \textit{\textbf{#1 \thepregun.}} % Inicio del bloque tipo cita
}{
    \end{quote} % Fin del bloque tipo cita
    \vspace{0.5em} % Espacio después del enunciado
}

\newenvironment{ejercicio}[1][]{\refstepcounter{ejer}\begin{quote}\textit{\textbf{#1 \theejer. }}}{\end{quote}\vspace{0em}}

% Define el entorno answer
\newenvironment{respuesta}{%
    \par % Mirar despues si queda bien el nuevo parrafo
    \noindent % El contenido en formato normal
    \ignorespaces
}{%
    \vspace{1.5em} % Espacio después de la respuesta
}

\newenvironment{solution}{\begin{proof}[Solution]}{\end{proof}}

\usepackage{pgfplots}
\usepackage{pgfplotstable}
\usepackage{booktabs} % for nicer table rules

% cargar datos de las mediciones
% cuadratico
\pgfplotstableread[col sep=space,header=true]{../data/burbuja.dat}\burbujadata
\pgfplotstableread[col sep=space,header=true]{../data/insercion.dat}\inserciondata
\pgfplotstableread[col sep=space,header=true]{../data/seleccion.dat}\selecciondata
\pgfplotstableread[col sep=space,header=true]{../data/shellsort.dat}\shellsortdata
\pgfplotstableread[col sep=space,header=true]{../data/shellsort_peor.dat}\shellsortpeordata
\pgfplotstableread[col sep=space,header=true]{../data/quicksort_peor.dat}\quicksortpeordata

% n log n
%\pgfplotstableread[col sep=space,header=true]{../data/mergesort.dat}\mergesortdata
%\pgfplotstableread[col sep=space,header=true]{../data/heapsort.dat}\heapsortdata
%\pgfplotstableread[col sep=space,header=true]{../data/heapsort_peor.dat}\heapsortpeordata
%\pgfplotstableread[col sep=space,header=true]{../data/quicksort.dat}\quicksortdata

% exponenciales
\pgfplotstableread[col sep=space,header=true]{../data/hanoi.dat}\hanoidata
\pgfplotstableread[col sep=space,header=true]{../data/fibonacci.dat}\fibonaccidata

\begin{document}
 
% --------------------------------------------------------------
%                         Start here
% --------------------------------------------------------------

% importar los valores de ajuste hibrido
\newcommand{\paramHanoi}{\num{4.71889e-08}}
\newcommand{\paramFibonacci}{\num{2.57762e-09}}

\newcommand{\paramInsercion}{\num{7.66757e-10}}
\newcommand{\paramSeleccion}{\num{1.21257e-09}}
\newcommand{\paramBurbuja}{\num{1.97198e-09}}
\newcommand{\paramShellsort}{\num{1.87424e-12}}
\newcommand{\paramQuicksortPeor}{\num{1.62547e-09}}

\newcommand{\paramMergesort}{\num{8.84509e-09}}
\newcommand{\paramMergesortPeor}{\num{3.01384e-08}}
\newcommand{\paramHeapsort}{\num{1.40579e-08}}
\newcommand{\paramQuicksort}{\num{1.11192e-08}}



\title{Práctica 1: cálculo de eficiencia de algoritmos}
\author{Juan Ignacio Molina Cobo \\ etc}

\date{}

\maketitle

\section{Introducción}

A continuación se plasman los resultados de elaborar la Práctica 1 de la asignatura.
Vamos a trabajar con los algoritmos propuestos (varios de ordenación, Hanoi, Fibonacci).
Para cada uno de ellos mostramos su eficiencia empírica a través de los tiempos de ejecución
de cada uno de ellos. Representaremos gráficamente estos datos y aplicaremos un modelo de ajuste
para calcular la eficiencia híbrida de cada uno de los algoritmos. El susodicho proceso se realizará
agrupando a los algoritmos en función de su eficiencia teórica. Adicionalmente, compararemos todos
los algoritmos de ordenación, independientemente de sus órdenes de eficiencia.

A la hora de agrupar, hemos diferenciado los casos promedio y peor de Quicksort. Mientras que el
primero lo hemos incluido dentro de los algoritmos $O(n \cdot\log n)$, el segundo lo hemos
includio dentro de $O(n^2)$.

Nota: esta memoria, así como los códigos fuente y scripts empleados y suficientes
para recrear los resultados aquí plasmados, es recuperable en https://www.github.com/igmc/p1-alg-eficiencia

\section{Cálculo de la eficiencia empírica}

\subparagraph{Algoritmos pertenecientes al orden $O(n^2)$}


Hemos elaborado un script shell para lanzar y recoger las mediciones para cada
uno de los algoritmos implementados. A continuación se muestran las tablas
mostrando, para cada orden de eficiencia, las mediciones obtenidas para cada 
algoritmo.

% Juntar los datos de las distintas tablas de O(n^2)

\pgfplotstablecopy{\burbujadata}\to{\cuadraticadata}

\pgfplotstablecreatecol[copy column from table={\inserciondata}{t_insercion}]
  {t_insercion}{\cuadraticadata}
\pgfplotstablecreatecol[copy column from table={\selecciondata}{t_seleccion}]
  {t_seleccion}{\cuadraticadata}
\pgfplotstablecreatecol[copy column from table={\shellsortdata}{t_shellsort}]
  {t_shellsort}{\cuadraticadata}
\pgfplotstablecreatecol[copy column from table={\shellsortpeordata}{t_shellsort_peor}]
  {t_shellsort_peor}{\cuadraticadata}
\pgfplotstablecreatecol[copy column from table={\quicksortpeordata}{t_quicksort_peor}]
  {t_quicksort_peor}{\cuadraticadata}

\sisetup{group-separator={}, group-minimum-digits=4}

% Mostrar tabla
\begin{table}[htbp]
  \centering
  \pgfplotstabletypeset[
    columns={n, t_burbuja, t_insercion, t_seleccion, t_shellsort, t_shellsort_peor, t_quicksort_peor},
  columns/n/.style={
    column name={$n$},
    assign cell content/.code={
      \pgfkeyssetvalue{/pgfplots/table/@cell content}{\num{##1}}
    }
  },
    columns/t_burbuja/.style={
      column name={Burbuja},
      fixed, fixed zerofill, precision=3
    },
    columns/t_insercion/.style={
      column name={Insercion},
      fixed, fixed zerofill, precision=3
    },
    columns/t_seleccion/.style={
      column name={Seleccion},
      fixed, fixed zerofill, precision=3
    },
    columns/t_shellsort/.style={
      column name={Shellsort},
      fixed, fixed zerofill, precision=3
    },
    columns/t_shellsort_peor/.style={
      column name={SS peor},
      fixed, fixed zerofill, precision=3
    },
    columns/t_quicksort_peor/.style={
      column name={QS peor},
      fixed, fixed zerofill, precision=3
    },
    every head row/.style={before row=\toprule, after row=\midrule},
    every last row/.style={after row=\bottomrule},
  ]{\cuadraticadata}
  \caption{tiempos de ejecución (s) de los algoritmos de orden $O(n^2).$}
  \label{tab:cuadratica}
\end{table}
% Juntar los datos de las distintas tablas

\pgfplotstablecreatecol[copy column from table={\hanoidata}{t_hanoi}]
  {t_hanoi}{\fibonaccidata}

\sisetup{group-separator={\ }, group-minimum-digits=4}

% Table
\pgfplotstabletypeset[
  columns={n, t_fibonacci, t_hanoi},
columns/n/.style={
  column name={$n$}
},
  columns/t_hanoi/.style={
    column name={Hanoi},
    precision=3
  },
  columns/t_fibonacci/.style={
    column name={Fibonacci},
    precision=3
  },
  every head row/.style={before row=\toprule, after row=\midrule},
  every last row/.style={after row=\bottomrule},
]{\fibonaccidata}
y ahora
 \paramHanoi
% --------------------------------------------------------------
%     You don't have to mess with anything below this line.
% --------------------------------------------------------------
 
\end{document}
